% LSPV-4 REPORT::CERTDPLL-BRIDGE1::2026-10-13T14:55:45+09:00
% !TEX program = xelatex
\documentclass[11pt]{article}
\usepackage[T1]{fontenc}
\usepackage{lmodern}
\usepackage[margin=0.86in]{geometry}
\usepackage{amsmath,amssymb,amsthm,booktabs,tabularx,xurl,hyperref}
\newtheorem{theorem}{Theorem}
\newtheorem{corollary}[theorem]{Corollary}
\newcommand{\Ltwo}{L_{22}}
\newcommand{\Res}{\mathsf{Res}}
\newcommand{\sCP}{\mathsf{sCP}}
\newcommand{\EXACT}{\textbf{PROVED/EXACT}}
\newcommand{\KNOWN}{\textbf{KNOWN THEOREM}}
\newcommand{\COMP}{\textbf{COMPUTATIONAL}}
\newcommand{\COND}{\textbf{CONDITIONAL}}
\newcommand{\CAND}{\textbf{CANDIDATE}}
\newcommand{\OPEN}{\textbf{OPEN}}
\newcommand{\RET}{\textbf{FALSIFIED/RETIRED}}
\hypersetup{colorlinks=true,urlcolor=blue,linkcolor=blue}
\setlength{\emergencystretch}{3em}
\title{CERTDPLL-BRIDGE1: A Candidate Restricted Machine-to-Proof Bridge\\\large Certificate-producing DPLL, not arbitrary one-bit computation}
\author{UNIVERSAL\_CORE\_BRIDGE / LSPV-4}
\date{13 October 2026, 14:55 KST}

\begin{document}
\maketitle

\begin{center}
\fbox{\parbox{0.94\linewidth}{
\textbf{Status: P versus NP remains OPEN.}
For a solver that writes an explicit DPLL derivation on every rejected input,
there is a fully charged, same-input translation to Resolution and hence to
semantic Cutting Planes.  The number of Resolution inferences does not exceed
the DPLL derivation size; a conservative explicit-bit implementation costs
only a further polynomial factor in the encoded input length.  This closes a
candidate restricted computation-to-proof arrow, with independent research
review still pending.  It does not apply to an
arbitrary total uniform decider that returns one bit.  The latter implication
remains OPEN and is not obtained from a computation trace.}}
\end{center}

\section{Claim card and fixed target}
Let $\mathcal C_{22}$ be the declared domain of nonempty proper
distinct-clause signed exact-$(2,2)$ 3-CNFs under the canonical binary
encoding, and let
\[
 \Ltwo=\{\operatorname{enc}(H):H\in\mathcal C_{22},\ H\text{ satisfiable}\},
\]
with malformed and out-of-domain strings rejected.

\begin{tabularx}{\textwidth}{@{}lX@{}}
\toprule
Claim ID & \textsf{CERTDPLL-TO-SCP1}.\\
Status before run & \CAND{}, isolated by RS5REVIEW1 as the weakest restricted transfer worth testing.\\
Desired conclusion & An explicit polynomial-time DPLL-certificate solver for $\Ltwo$ yields polynomial-size/depth $\sCP$ refutations on the same rejected inputs.\\
Exact prefix & $\forall A\,\forall H\in\mathcal C_{22}$: if $A(H)$ rejects and its same run writes a valid DPLL derivation $\Pi_H$, construct a $\Res$ and $\sCP$ refutation with all reading, checking, translation, output and workspace charged.\\
Machine/model & Uniform deterministic bit machine; explicit self-delimiting proof output; runtime includes every output bit.\\
SAME\_REALIZATION & The encoded formula $H$, run $A(H)$, emitted $\Pi_H$, translated $\Res$ proof and $\sCP$ proof all refer to the same $H$.\\
What would falsify it & A valid DPLL derivation whose standard translation is larger than polynomial in the explicit input/certificate length, or whose clauses are not semantic-CP lines.\\
\bottomrule
\end{tabularx}

The unrestricted project goal is still
\[
 \forall A\ \forall k\in\mathbb N\ \forall C>0\ \forall N_0\ \exists y:
 |y|\ge N_0,\qquad T_A(y)>C|y|^k,                         \tag{1}
\]
where $A$ ranges over all correct total deterministic uniform deciders of
$\Ltwo$.  The theorem below concerns a strict subclass and does not prove
(1).

\section{Primary-literature preflight}
The two obstacles checked first were proof emission versus one-bit decision,
and whether the DPLL-to-Resolution translation has a quantitative bound.

Berger, Lawrence, Nordvall Forsberg and Seisenberger define DPLL derivations
with Unit, Red, Elim, Conflict and Split rules and prove their Theorem~6:
for every consistent valuation $\Gamma$, if a DPLL derivation of
$\Gamma\vdash\Delta$ has size $n$, then $\Delta$ has a Resolution derivation
of $\overline\Gamma$ of size $m\le n$ \cite{berger}.  For empty $\Gamma$ this
is a Resolution refutation.  This is a \emph{direct theorem transfer} to an
explicit DPLL certificate.

Buss, Hoffmann and Johannsen characterize DLL algorithms with clause learning
by named Resolution refinements and make their simulation assumptions
explicit \cite{buss}.  Beame, Kautz and Sabharwal likewise characterize
specified clause-learning systems, while emphasizing nondeterministic proof
search choices \cite{beame}.  These are direct transfers only for the
declared solver/proof models, not for arbitrary SAT programs.

The current source size--depth theorem is about Resolution, semantic Cutting
Planes and triangle DAGs in their named models \cite{tradeoff}.  Treating any
one-bit execution as such a proof is a \emph{failed transfer}, already exposed
by TRACEEXT-FIREWALL1.  Circuit-observability, VLSI tracing and cryptographic
special-soundness remain analogies requiring a new bridge.

\section{Candidate restricted bridge}
Call a total solver $A$ \emph{explicit CERT-DPLL} if, on a satisfiable CNF, it
writes a satisfying assignment, and on an unsatisfiable CNF it writes a
self-delimiting DPLL derivation using the rules in \cite{berger}.  Let $N$ be
the bit length of $\operatorname{enc}(H)$ and let $T_A(H)$ count the complete
run, including certificate output.

\begin{theorem}[certificate-producing DPLL to Resolution and $\sCP$]
Suppose $H$ is unsatisfiable and an explicit CERT-DPLL solver emits a valid
DPLL derivation $\Pi_H$ with $n$ rule occurrences and split depth $d$.  Then:
\begin{enumerate}
\item $H$ has a Resolution refutation with at most $n$ Resolution inferences;
\item it has a semantic Cutting Planes refutation with at most $n+O(|H|)$
lines, depth at most $n+O(1)$ and coefficients in $\{-1,0,1\}$;
\item since the certificate is explicit, $n\le |\operatorname{enc}(\Pi_H)|
\le T_A(H)$ up to the fixed machine's output convention;
\item a conservative materializing translator runs in
$O(T_A(H)N\log(N+2))$ bit time and uses at most that many output/work bits.
All objects refer to the same formula $H$.
\end{enumerate}
\end{theorem}

\begin{proof}
Apply Theorem~6 of \cite{berger} with empty initial valuation.  It produces a
Resolution derivation of the empty clause with $m\le n$ inferences.  For the
$\sCP$ embedding, represent a clause $C$ by
\[
 \sum_{x\in C}x+\sum_{\neg x\in C}(1-x)\ge1.             \tag{2}
\]
Every Boolean assignment satisfying both parent clauses of a Resolution
inference satisfies its resolvent, so the corresponding line is a legal
semantic inference.  Axiom clauses of $H$ and Boolean bounds supply the
initial lines.  The coefficients in (2) have constant bit length.

An explicit output requires at least one written symbol per rule occurrence,
so $n$ is bounded by output length and hence runtime.  A simple translator
may materialize each derived clause as a sorted list of at most $N$ literals,
with indices of $O(\log(N+2))$ bits.  Reading, checking, translating and
writing all at most $T_A(H)$ rule occurrences gives the displayed conservative
bound.  This charges proof acquisition and output instead of treating proof
existence as free.
\end{proof}

\begin{corollary}[restricted lower-bound consequence]
Any proved $\sCP$ size--depth lower bound rules out, on the same family, an
explicit CERT-DPLL implementation whose complete runtime would force both
translated size and depth below that lower-bound region.  DIAGTRANSFER1's
exponent-indexed selection may be applied after fixing the solver exponent.
This corollary is restricted to proof-emitting solvers.
\end{corollary}

\paragraph{Sharper split-only form.}
For the elementary split tree used in the bounded audit, each leaf at path
valuation $\Gamma$ chooses an input clause $C\subseteq\overline\Gamma$ and
uses subsumption to derive $\overline\Gamma$.  An internal split on $x$
resolves
$\overline\Gamma\vee\neg x$ with
$\overline\Gamma\vee x$ to obtain $\overline\Gamma$.  Hence the number of
Resolution inferences equals the number of split nodes and Resolution depth
equals split depth.  This is a direct symbolic proof, not an inference from
finite tests.

\section{Why this does not cover arbitrary deciders}
An arbitrary decider for $\Ltwo$ is required to output only one bit.  Neither
its specification nor polynomial runtime requires it to emit a DPLL
derivation, a Resolution proof, a semantic inequality, or a triangle
relation.  A wrapper may compute by Gaussian elimination, compact arithmetic,
dynamic programming, a proof system of different strength, or any other
representation and discard all auxiliary data before returning the bit.

TRACEEXT-FIREWALL1 proves a sharper obstruction to the naive repair: the
Tseitin encoding of a fixed evaluated run plus the opposite output has a
linear unit-Resolution refutation, but substituting actual trace values
satisfies every computation axiom and leaves only the deliberately false
output unit.  It is not a proof from the clauses of $H$.  Thus
\[
 \boxed{
 \text{explicit DPLL certificate}\Rightarrow\Res\Rightarrow\sCP
 }
 \quad\text{is \CAND{} pending independent review,}        \tag{3}
\]
whereas
\[
 \boxed{
 \text{arbitrary one-bit decider}\Rightarrow
 \text{explicit DPLL/$\sCP$ certificate}
 }
 \quad\text{remains \OPEN{}.}                              \tag{4}
\]
Promoting (4) would already exclude every polynomial-time decider on the hard
family and is therefore not a harmless implementation lemma.

\section{Bounded reproducible audit}
The independent finite checker
\path{experiments/review/dpll_certificate_bridge_audit.py} uses seed
\texttt{202610131455}.  It generated explicit split trees, translated leaves
by subsumption and internal nodes by Resolution, then checked every source
inclusion, pivot, resolvent, root empty clause, proof size and depth.

In 0.089264 seconds it completed 19,877 assertions over 654 formulas:
263 satisfiable controls and 391 unsatisfiable formulas.  The unsatisfiable
proofs contained 3,989 proof nodes and 1,799 Resolution inferences.  The audit
included Horn, 2-SAT/Krom, affine/XOR, bounded-treewidth, independent-product,
renaming, phase, clause-order and tautology-padding controls, plus 240
DEAD/COPY/DELAY/INTERP/RECOMP certificate-preservation checks.  It separately
confirmed that bare output bits are not accepted as certificates.

Raw output is
\path{results/raw/review/2026-10-13_1455_dpll_certificate_bridge.json}.
SHA-256 values are
\texttt{e0b756833bdf1945de8fc6d02c8205586d2ba93d1fb601f20d57c047c2d25873}
for the script and
\texttt{6d3651c6ab099682d5e500d0c0d47139750a11092f8aec44685642406af129f3}
for the result.  The semantic digest inside the result is
\texttt{92494b0ac39b89d8125585d16a10a24e2536663f41917f7ec314d9ffa1ed65f2}.
The finite audit is \COMP{} only.

\section{Hostile controls and barrier audit}
\paragraph{Easy families and alternate representations.}
The bridge behaves correctly on Horn, Krom, affine and bounded-treewidth
inputs because it charges only an explicitly emitted certificate.  It does
not claim these families are intrinsically hard.  Gaussian elimination and
compact DP are genuine alternate implementations outside CERT-DPLL and show
why (4) cannot be inferred from (3).

\paragraph{Wrappers.}
DEAD, COPY, DELAY, INTERP and RECOMP may increase runtime while preserving the
same emitted certificate, so they do not falsify (3).  A wrapper that erases
the certificate and emits one bit leaves the CERT-DPLL class and lands exactly
at the open arrow (4).  This is a specification boundary, not an invariant of
arbitrary programs.

\paragraph{Relativization and algebrization.}
The syntactic translation of a supplied certificate survives oracle access,
but an oracle decider need not supply a certificate.  Hence no oracle
separation follows.  Low-degree or algebraic access may make affine controls
easier without affecting the conditional translation.  No nonrelativizing or
nonalgebrizing hinge is obtained.

\paragraph{Natural Proofs and post-quantum ledger.}
No large constructive circuit property is defined.  Schur squares,
quadratic-relation ranks, Hilbert functions, syzygies and collision entropy
are absent from the certificate translation and are \textbf{FAILED TRANSFER}
for this barrier.  LWE-style external extraction and special-soundness remain
\textbf{ANALOGY/CANDIDATE}: they would need an unconditional black-box theorem
producing (3)'s proof object from a one-bit solver.

\section{Corpus reconciliation}
\paragraph{CORPUS REVIEWED.}
The run surveyed the current GitHub tree and read root \path{main.tex},
\path{AGENTS.md}, \path{CODEX.md}, the complete pnp-research skill and its mode
and protocol references, every current task in \path{tasks/}, actual
\path{main4.tex} and \path{main5.tex}, newest MAIN RS5REVIEW1, DIAGTRANSFER1,
ROUNDINGSQUEEZE5, latest LSPV-1 through LSPV-4 reports, the proof-agent
REDTEAM-STRENGTHEN12 review, both current governance memos, TRACECR-FIREWALL1,
TRACEAXIOM-FIREWALL1, IMAGEEQUIV1, ARC22RETURN1, SOURCEARC-NOCONTRACT1,
PC22\_ALLFIELD1, and the HEC/VSC-HEC/G-PPC-HEC/CRS-PPC/CCSPC causal branch.
The immediately preceding LSPV-3 COSPECBLOWRED1 report is substantive and has
its required marker and sentinel; no predecessor recovery was due.

The exact local body of the previously unpublished LSPV-4 report
TRACEEXT-FIREWALL1 was recovered without reconstruction to
\path{research/agents/lspv-4/2026-10-13_1402_LSPV-4_TRACEEXT-FIREWALL1.tex}
and verified before this report.

\paragraph{DUPLICATION / CONTRADICTION AUDIT.}
The DPLL-to-Resolution simulation is a known theorem and is cited rather than
claimed as new.  TRACESIMQ1, TRACECR-FIREWALL1 and TRACEAXIOM-FIREWALL1 already
separate decider-specific proof systems and fixed traces from preselected
lower-bounded systems.  The new project contribution is the candidate positive
boundary (3), with complete resource accounting and a tested explicit
certificate format.  It does not contradict TRACEEXT-FIREWALL1 because a
DPLL proof has clauses of $H$ as premises, whereas the fixed execution trace
has computation axioms and a false output premise.

\paragraph{WHAT CHANGED RELATIVE TO ROOT MAIN.}
Root main correctly records the unrestricted bridge as OPEN.  This run adds
one positive restricted entry: explicit CERT-DPLL computations translate to
Resolution/$\sCP$ with polynomial charged overhead.  The executive proof
spine and P-versus-NP status do not change.

\paragraph{ADOPTED / REJECTED / RETIRED PRIOR COMPONENTS.}
Adopt RS5REVIEW1's independently verified triangle-DAG lifting theorem,
DIAGTRANSFER1's exponent-indexed composition, IMAGEEQUIV1's exact image
equivalence, and the existing compiler/projection results only within their
declared models.  Preserve ARC22RETURN1's typed depadding and expanding
restriction--reencoding result.  Reject any universal inference from (3).
Retire the claim that proof logging can be added to an arbitrary decider with
polynomial overhead without proving a new theorem.

\paragraph{STALE OR OVERSTATED CLAIMS FOUND.}
The blanket phrase ``no machine-to-proof transfer is known'' is now too broad.
A standard transfer is known for explicit DPLL-certificate solvers.  The
correct unresolved statement is that no transfer is known for every total
uniform one-bit decider of $\Ltwo$ into the lower-bounded $\sCP$/triangle
model.

\section{Integration acknowledgments}
\paragraph{EARLY22-INTEGRATION1.}
The exact-$(2,2)$ compiler/minimality/Resolution projection results remain
preserved at their graph, pairing and proof-system hypotheses, separately from
the monotone cubic-pair theorem and affine-continuation limitation.  Historical
$N$ is a variable count; encoded length remains $\Theta(N\log(N+2))$ where
applicable.  Exact manuscript location is this section and the LSPV-4 owned
region of root \path{main.tex}.  Next check: replay the full selector-compiler
certificate batch before promoting the complete proof-model composition.

\paragraph{SINGLE22-SEPARATION1.}
The target remains the one fixed NP-complete language $\Ltwo$.  This run
advances exactly one arrow: proof-emitting DPLL computation to the named
lower-bounded proof model.  It does not require a broader SAT family and does
not confuse one formula or a NO benchmark with the language.  The additional
arrow needed for (1) is still (4).

\section{Status and next target}
\begin{tabularx}{\textwidth}{@{}l l X@{}}
\toprule
Component & Status & Reason\\
\midrule
DPLL derivation of size $n$ to Resolution size at most $n$ & \KNOWN & Berger et al., Theorem~6.\\
Explicit CERT-DPLL to charged $\Res/\sCP$ & \CAND & Direct composition and clause-inequality embedding; independent research review pending.\\
Finite certificate translator audit & \COMP & 19,877 assertions; no asymptotic use.\\
Fixed evaluated trace as input proof & \RET & TRACEEXT-FIREWALL1.\\
Arbitrary one-bit decider to CERT-DPLL/$\sCP$ & \OPEN & No theorem forces proof emission or fixed-system simulation.\\
Whole-computation lower bound (1) & \OPEN & The universal transfer is load-bearing.\\
\bottomrule
\end{tabularx}

\paragraph{Highest-value next target.}
Test a broader but still exact proof-logging class: CDCL with a declared
learning rule, restarts and explicit proof log.  Use the primary simulation
theorems to record exact size/depth and output costs, then stop at the first
rule whose translation is not polynomial.  This can enlarge the restricted
positive region without pretending to cover arbitrary deciders.

P versus NP remains open.

\begin{thebibliography}{9}
\bibitem{berger}
U. Berger, A. Lawrence, F. Nordvall Forsberg, and M. Seisenberger,
``Extracting Verified Decision Procedures: DPLL and Resolution,''
\emph{Logical Methods in Computer Science} 11(1:6), 2015.
\url{https://arxiv.org/abs/1502.02131}.

\bibitem{buss}
S. R. Buss, J. Hoffmann, and J. Johannsen,
``Resolution Trees with Lemmas: Resolution Refinements that Characterize DLL
Algorithms with Clause Learning,'' \emph{Logical Methods in Computer Science}
4(4), 2008.
\url{https://arxiv.org/abs/0811.1075}.

\bibitem{beame}
P. Beame, H. Kautz, and A. Sabharwal,
``Towards Understanding and Harnessing the Potential of Clause Learning,''
\emph{Journal of Artificial Intelligence Research} 22, 2004.
\url{https://arxiv.org/abs/1107.0044}.

\bibitem{tradeoff}
S. F. de Rezende, N. Fleming, D. A. Janett, J. Nordstr\"om, and S. Pang,
``Truly Supercritical Trade-Offs for Resolution, Cutting Planes, Monotone
Circuits, and Weisfeiler--Leman,'' arXiv:2411.14267; STOC 2025.
\url{https://arxiv.org/abs/2411.14267}.
\end{thebibliography}

\end{document}
% END_OF_VERIFIED_REPORT::LSPV-4::CERTDPLL-BRIDGE1
